\section{Diseño del Algoritmo}

. explico la logica del algoritmo y como llegamos a eso
. muestro el codigo (entero? entra/me sirve? si es corto ponerlo, y si no vemos.)
. explico el por qué hago la implementacion de esa manera

\subsection{Lógica}

\subsection{Código}




\begin{lstlisting}[language=Python]
import sys

PRIMERO = 0
ULTIMO = -1

def juego(monedas: list[int]):
    pts_mateo = pts_sophia = 0
    
    turno = True
    
    while len(monedas) > 0:
        if turno:
            if monedas[PRIMERO] > monedas[ULTIMO]:
                pts_sophia += monedas[PRIMERO]
                monedas[:] = monedas[PRIMERO+1:]
                sys.stdout.write("Primera moneda para Sophia;")
            else:
                pts_sophia += monedas[ULTIMO]
                monedas[:] = monedas[:ULTIMO]
                sys.stdout.write("Ultima moneda para Sophia;")
        else:
            if monedas[PRIMERO] > monedas[ULTIMO]:
                pts_mateo += monedas[ULTIMO]
                monedas[:] = monedas[:ULTIMO]
                sys.stdout.write("Ultima moneda para Mateo;")
            else:
                pts_mateo += monedas[PRIMERO]
                monedas[:] = monedas[PRIMERO+1:]
                sys.stdout.write("Primera moneda para Mateo;")
        turno = not turno
        
\end{lstlisting}
